% Testkapitel für die Pipeline. Eigener Text, kein Buchtext.
\chapter[Testkapitel]{}
\kapitelauftakt{9}{Ein Testkapitel}

\section{Erster Abschnitt}
\label{sec:k9-erster}

\grossdiagrammtaktik[\abbziel{9.1}: Nach 1. e4 -- \textbf{Ziel}: das Feld \glqq e4\grqq{} (gold).]
{rnbqkbnr/pppppppp/8/8/4P3/8/PPPP1PPP/RNBQKBNR b KQkq e3 0 1}
{gold/e2-e4}
{gold/e4,grau/d5}

\subsection*{Eine Reihe}

\begin{diagrammreihe}
\diagrammschritttaktik
{rnbqkbnr/pppp1ppp/8/4p3/4P3/8/PPPP1PPP/RNBQKBNR w KQkq e6 0 2}
{rot/e5-e4}
{}
{\abbziel{9.2}: Nach 1... e5, Weiß am Zug. Siehe \abb{9.1} und \abbnr{9.3}.}
\diagrammschritt{rnbqkbnr/pppp1ppp/8/4p3/4P3/5N2/PPPP1PPP/RNBQKB1R b KQkq - 1 2}{\abbziel{9.3}: Nach 2. Sf3 droht \zugfolge{Sxe5}.}
\end{diagrammreihe}

\section{Zweiter Abschnitt}

\begin{vorsichtfalle}
\diagrammbox[\abbziel{9.4}: Schwarz am Zug, $-$3 für Weiß.]{rnbqkbnr/pppp1ppp/8/4p3/4P3/5N2/PPPP1PPP/RNBQKB1R w KQkq - 1 2}
\end{vorsichtfalle}
